\documentclass[11pt]{article}
\newcommand{\ShortTitle}{HODGE / CYCLES, PACKETS AND CARRIERS}
\usepackage[T1]{fontenc}
\usepackage[utf8]{inputenc}
\usepackage{mathpazo}
\usepackage[scaled=0.92]{helvet}
\usepackage{courier}
\usepackage{amsmath,amssymb,mathtools}
\usepackage[a4paper,landscape,left=18mm,right=18mm,top=19mm,bottom=17mm,headheight=14pt]{geometry}
\usepackage{graphicx,microtype,xcolor,fancyhdr,hyperref,booktabs,array}
\definecolor{Ink}{HTML}{243138}
\definecolor{Teal}{HTML}{21676B}
\definecolor{Copper}{HTML}{A45C36}
\definecolor{Soft}{HTML}{E9F0EF}
\definecolor{Gray}{HTML}{627078}
\hypersetup{colorlinks=true,linkcolor=Teal,citecolor=Teal,urlcolor=Teal,pdfauthor={Komorebi P.}}
\urlstyle{same}
\setlength{\parindent}{0pt}
\setlength{\parskip}{0.55em}
\setlength{\emergencystretch}{1.5em}
\linespread{1.04}
\pagestyle{fancy}
\fancyhf{}
\fancyhead[L]{\sffamily\scriptsize\color{Gray}OPERATION KOMOREBI / VISUAL GUIDE}
\fancyhead[R]{\sffamily\scriptsize\color{Gray}\ShortTitle}
\fancyfoot[L]{\sffamily\scriptsize\color{Gray}READER EDITION 1.0 / 9 OCTOBER 2026}
\fancyfoot[R]{\sffamily\small\color{Ink}\thepage}
\renewcommand{\headrulewidth}{0.25pt}
\renewcommand{\footrulewidth}{0pt}
\providecommand{\N}{\mathbb N}
\providecommand{\Z}{\mathbb Z}
\providecommand{\Q}{\mathbb Q}
\providecommand{\R}{\mathbb R}
\providecommand{\C}{\mathbb C}
\providecommand{\eps}{\varepsilon}
\providecommand{\dd}{\mathrm d}
\providecommand{\norm}[1]{\left\lVert #1\right\rVert}
\providecommand{\abs}[1]{\left\lvert #1\right\rvert}
\newcommand{\GuideCover}[4]{%
\hypersetup{pageanchor=false}
\begin{titlepage}\thispagestyle{empty}
{\sffamily\small\color{Teal}OPERATION KOMOREBI\hfill CONJECTURE SERIES #1 / VISUAL GUIDE}\par
\vspace{18mm}
{\fontsize{40}{44}\selectfont\color{Ink}#2}\par
\vspace{4mm}{\LARGE\color{Ink}#3}\par
\vspace{6mm}{\color{Teal}\rule{\linewidth}{1.1pt}}\par
\vspace{4mm}{\large Komorebi P.}\par
{\sffamily\small\color{Gray}READER EDITION 1.0 / 9 OCTOBER 2026}\par
\vfill
\begin{minipage}{0.78\linewidth}\large #4\end{minipage}\par
\vspace{5mm}{\sffamily\small\color{Teal}WORKED EXAMPLES\quad /\quad EXACT COORDINATES\quad /\quad OPEN PROBLEMS}\par
\end{titlepage}\hypersetup{pageanchor=true}}
\newcommand{\GuideSpread}[5]{%
\clearpage
{\fontsize{22}{26}\selectfont\bfseries\color{Ink}#1}\par
\vspace{1mm}{\sffamily\small\color{Gray}#2}\par
\vspace{3mm}
\noindent\begin{minipage}[t]{0.665\linewidth}\vspace{0pt}
\centering\includegraphics[width=\linewidth,height=113mm,keepaspectratio]{#3}
\end{minipage}\hfill
\begin{minipage}[t]{0.295\linewidth}\vspace{0pt}
{\sffamily\small\bfseries\color{Teal}READ THE FIGURE}\par
\small #4
\end{minipage}\par
\vfill
\noindent\colorbox{Soft}{\parbox{\dimexpr\linewidth-2\fboxsep\relax}{\vspace{1.3mm}\small\textbf{Carry forward.} #5\vspace{1.3mm}}}\par
}

\hypersetup{pdftitle={Hodge Conjecture: Visual Guide},pdfsubject={Exact visual guide to rational cycles, Fermat packets, constructive carriers and remaining gaps}}
\graphicspath{{../shared/figures/}}
\begin{document}
\GuideCover{06}{The Hodge Conjecture}{A visual guide to cycles, packets and carriers}{The conjecture asks whether a rational cohomology class of the right Hodge type comes from algebraic cycles. These spreads follow the precise gap, the finite certificates that can close a packet, the constructions that work, and the obstructions that force a change of method. The general conjecture remains open.}

\GuideSpread{01 / The gap is between two rational spaces}{SCHEMATIC / THE MATHEMATICAL TARGET}{01_hodge_gap.pdf}{Every algebraic codimension-$p$ cycle has a rational class of type $(p,p)$. This gives the inner space. The outer space contains all rational cohomology classes of that type.\par
The question mark denotes a possible unrepresented direction. Its size is not a measurement, and the drawing does not claim that a gap exists.\par
The conjecture asserts equality for every smooth projective complex variety and every $p$.}{A type condition identifies a candidate class. A proof must supply its expression in algebraic cycle classes. Primary statement: Deligne [1].}

\GuideSpread{02 / Rational combinations are the target}{SCHEMATIC / A COEFFICIENT EXAMPLE}{02_rational_cycle_combination.pdf}{The cycle map sends a subvariety to its cohomology class. The target allows finite rational combinations of those classes.\par
The coefficients shown are an illustrative definition. A negative coefficient is permitted. A common denominator turns the expression into a rational multiple of an integral cycle.\par
This does not ask for a single effective surface with the prescribed class. The integral Hodge conjecture is a different statement and is false in general.}{Keep rationality visible throughout the construction. A complex eigenvector with the right type is not yet a rational cycle class.}

\GuideSpread{03 / Why a fourfold focuses attention on the middle}{EXACT INDEXING / HODGE BIDEGREES}{03_fourfold_middle.pdf}{For complex dimension four, $0\le p,q\le4$. The row with $p+q=4$ is the middle cohomology. Its central position is $(2,2)$, the type of a surface class.\par
The cells show possible bidegrees. They are not numerical Hodge dimensions.\par
The work here studies primitive middle packets on Fermat fourfolds. This is a specific testing ground, with additional diagonal symmetry and explicit residue calculations.}{A construction on a special fourfold does not reduce arbitrary smooth projective varieties to that example.}

\GuideSpread{04 / An obstruction projector diagnoses the question}{SCHEMATIC / A REFORMULATION}{04_projector_diagnostic.pdf}{Extend the rational spaces to real vector spaces. Choose a positive inner product and take orthogonal projectors onto their real extensions. Since the latter lies in the former, their difference detects the unrepresented directions.\par
Writing $O=0$ restates equality of the spaces.\par
Neither naming $P_{\mathrm{Alg}}$ nor diagonalizing a Hodge operator constructs the unknown cycles. An algebraic correspondence representing an operator needs a separate geometric argument.}{Projector language organizes the task. It does not remove the cycle-existence obligation.}

\GuideSpread{05 / Positivity can give the wrong answer}{EXACT RATIONAL COUNTEREXAMPLE / CORRECTION TO THE OLD FRAME CLAIM}{05_frame_counterexample.pdf}{Take $V=\Q e_1$ and one candidate class $c=e_1+e_2$. The frame is $S=cc^{\mathsf T}$. Its compression has $\langle e_1,Se_1\rangle=1>0$.\par
Yet $e_1$ does not lie in the line spanned by $c$. The augmented matrix has rank two, while $c$ alone has rank one.\par
Positivity of the compression to $V$ therefore cannot certify that cycles span $V$. The candidate class must already lie in $V$, or a valid invariance argument must be supplied.}{Use an actual algebraic packet projector on already algebraic cycles when one is available. Then test the resulting contained frame.}

\GuideSpread{06 / A contained frame gives a rational certificate}{EXACT TOY EXAMPLE / ALL THREE COLUMNS LIE IN THE PACKET}{06_corrected_rational_frame.pdf}{The three columns are $c_1=(1,0)$, $c_2=(0,1)$ and $c_3=(1,1)$. They lie in $V=\Q^2$.\par
The frame $S=CC^{\mathsf T}$ has determinant three, so its image is all of $V$. For $\alpha=(2,3)$, exact reconstruction gives
\[\lambda=C^{\mathsf T}S^{-1}\alpha=(1/3,4/3,5/3).\]
Thus $C\lambda=\alpha$ over $\Q$.\par
The arithmetic is a toy model. In geometry the columns must be verified algebraic cycle classes.}{For a general cycle matrix $C$ and target-basis matrix $H$, the safe membership test is $\operatorname{rank}_{\Q}C=\operatorname{rank}_{\Q}[C\ H]$.}

\GuideSpread{07 / Positive children need not be orthogonal}{EXACT TOY EXAMPLE / POLYPHASE ASSEMBLY}{07_positive_polyphase.pdf}{Partitioning the cycle indices partitions the rank-one frame terms. This yields a sum of positive operators.\par
Here $S_0$ sees the first coordinate, and $S_1$ sees the diagonal direction. Their image directions have nonzero inner product. Their kernels have zero intersection, so the parent is invertible.\par
For degree residues the exact identity is
\[S_{M,r}=\sum_{j=0}^{L-1}S_{LM,r+jM}.\]
It follows from a disjoint partition of indices, without assuming orthogonal geometric packets.}{The parent kernel equals the intersection of positive child kernels. A complete partition preserves the frame even when the subspaces overlap.}

\GuideSpread{08 / Support must be verified before multiplicities}{SCHEMATIC / COHOMOLOGY WITH SUPPORTS}{08_support_verification.pdf}{Exactness says that a class comes from cohomology supported on $Z$ precisely when restriction to $X\setminus Z$ vanishes.\par
This is a test of a proposed support. It does not find $Z$.\par
When the codimension-$p$ components generate the relevant supported classes, purity supplies their cycle classes. Intersections with test classes give a rational system $Mq=b$ for the coefficients.\par
The codimension and purity hypotheses are part of the result.}{A guessed support and a solved linear system do not replace the vanishing check on the complement.}

\GuideSpread{09 / Correspondences shift degrees by a fixed rule}{EXACT INDEXING / DIVISOR LOWERING}{09_correspondence_degrees.pdf}{Let $\dim X=n$ and $\dim Y=d$. A codimension-$r$ correspondence from $X$ to $Y$ shifts cohomological degree by $2(r-n)$.\par
To lower $H^{2p}(X)$ to $H^2(Y)$, choose $r=n-p+1$. A return correspondence has codimension $d+p-1$.\par
Rational Hodge classes of type $(1,1)$ on $Y$ are divisor classes by Lefschetz. Returning those divisors gives algebraic classes on $X$.}{The composite must preserve the chosen packet and be invertible there. Algebraicity of lowering and return maps is a geometric hypothesis.}

\GuideSpread{10 / The inverse can stay algebraic}{EXACT TOY EXAMPLE / CAYLEY--HAMILTON RECONSTRUCTION}{10_rational_inverse.pdf}{For $T=\left(\begin{smallmatrix}2&1\\1&2\end{smallmatrix}\right)$, the characteristic polynomial is $t^2-4t+3$. Hence the inverse is a rational polynomial in $T$.\par
If $T$ comes from algebraic correspondences, its powers and rational sums also act through algebraic correspondences.\par
This lets one recover $\alpha$ from an already algebraic $T\alpha$. The same argument works in any finite dimension when the constant characteristic-polynomial coefficient is nonzero.}{A rational linear inverse alone is insufficient. The polynomial formula ties the inverse to the same algebraic operator.}

\GuideSpread{11 / Balance must hold across the rational orbit}{EXACT DATA / THE DEGREE-70 CARRIER PACKET}{11_balanced_orbit70.pdf}{Each bar stacks the six least-positive residues of $u a$ modulo 70, with
\[a=(1,20,24,42,61,62).\]
All 24 unit multipliers give total $210=3m$. Each ordered conjugate is distinct.\par
This puts the entire rational packet in middle Hodge type $(2,2)$. Checking only the sum of $a$ would test one complex line.\par
The plot is a direct integer enumeration, independently replayed in the new edition.}{For a full-support Fermat fourfold packet, the balance equation is $\sum_i\langle ua_i\rangle_m=3m$ for every unit $u$. Source: carrier [C].}

\GuideSpread{12 / Degree and conductor are different addresses}{EXACT INTEGER DATA / THE SOURCE AND TARGET AT DEGREE 210}{12_conductor_descent.pdf}{The source is $s=3a$ at modulus 210. All entries share the factor three with the modulus, so its exact conductor is 70.\par
The target has gcd one with 210. Its conductor is therefore 210.\par
Both sit on the same degree-210 Fermat fourfold, but the target is not a coordinate-power inflation of a lower-conductor packet.\par
The formula is $m/\gcd(m,a_0,\ldots,a_5)$.}{Pullback produces an infinite tower from a solved seed. A full-conductor target requires an additional construction. Source: transfer [T], Sections 4--5.}

\GuideSpread{13 / A dyadic star has an explicit third-order saturation}{EXACT MULTISET CERTIFICATE / MODULUS 32}{13_dyadic_sat3.pdf}{The packet $A_{32}$ is the first full-conductor four-star representative. Adding three complementary pairs gives the three four-entry blocks on the right.\par
The multiplicity of every residue agrees exactly. Each four-entry block is balanced for all units modulo 32, so belongs to the surface branch.\par
SAT(3) here denotes this displayed three-block stabilized identity. It is a precise certificate, rather than a numerical approximation.}{The displayed SAT(3) identity is exact. The proposed all-degree classification remains a write-up obligation, separate from fresh exhaustive replay through degree 128. Source: August [A] and its later status correction.}

\GuideSpread{14 / The short defect becomes a tree boundary}{EXACT COMBINATORIAL DIAGRAM / THE INTERNAL TREE AT MODULUS 32}{14_dyadic_tree32.pdf}{The root is $q/2=4$ for $q=8$. The parent rule doubles toward the root, with the reflected branch $p(b)=2(q-b)$.\par
The highlighted vertices form a star with center two. Completing the tree by zero-valued boundary leaves makes each internal vertex trivalent.\par
A connected four-vertex support has boundary size $4+2=6$. That is the defect mass allowed by a six-entry packet.\par
Boundary leaves are omitted in the drawing, not in the norm formula.}{Tree total variation replaces a residue-lattice problem by an exact combinatorial one. The proposed all-degree classification requires a completed universal write-up. The finite census and displayed identities remain independently exact.}

\GuideSpread{15 / Degree 33 forced a cross-family route}{MANUSCRIPT GEOMETRY / FINITE LATTICE IDENTITIES}{15_klein33_route.pdf}{The exceptional packet escaped the earlier pair and standard grammar. A proposed symmetric surface gave only the ambient class $242h^2$, so its primitive projection vanished.\par
The useful route goes to a smooth Klein--Delsarte cubic. Its monomial lattice map has the displayed Smith diagonal, giving generic torus degree $33^4$.\par
A checked residue numerator matches the exceptional Fermat line. Cubic divisor lowering and a rational graph correspondence then provide the packet argument.}{Failure of one cycle family can reveal a missing correspondence. It does not show nonalgebraicity. Source: [A], degree-33 chapter and verifier.}

\GuideSpread{16 / K23 separates into curve, surface and curve}{EXACT MULTISET CERTIFICATE / MODULUS 50}{16_k23_factorization50.pdf}{The five nonzero mod-five phases split as a complementary pair and a zero-sum triple. This is the K23 pattern.\par
Adding the two displayed Tate pairs gives a three-entry curve packet, a four-entry surface packet and a lower curve packet.\par
This identity organizes the tensor factors. It does not by itself construct the Chow correspondence between the required factors.\par
The sample is regenerated from the first stored K23 certificate.}{Keep the exact multiset identity separate from the geometric morphism it suggests. Source: [A], K23 audit, first selected certificate.}

\GuideSpread{17 / The imaginary selector reads an exact Hodge sign}{EXACT SIGN DATA / A LOWER CURVE PACKET AT MODULUS 7}{17_cm_imaginary_sign.pdf}{For $\ell=(1,2,4)$ modulo seven, conjugation sends $\Delta_\ell$ to $-\Delta_\ell$.\par
At each embedding, the age
\[h_\ell(u)=\frac{\langle u\rangle_7+\langle2u\rangle_7+\langle4u\rangle_7}{7}\]
is one or two. The sign of the imaginary part is $(-1)^{h_\ell(u)}$.\par
The bars show signs derived from integer sums, with no numerical evaluation of roots of unity.}{A rational character idempotent selects the whole Galois orbit. The upper half-plane is an embedding-sector description, not a rational half-orbit projector.}

\GuideSpread{18 / The missing arrow is a Chow correspondence}{SCHEMATIC / AN OPEN CONSTRUCTION OBLIGATION}{18_missing_chow_bridge.pdf}{The candidate lower Fermat-curve factor and the weight-one tensor packet have compatible Hodge data.\par
A geometric proof still needs an algebraic correspondence with the prescribed character covariance, and a nonzero coefficient after character projection.\par
Rank one makes a nonzero morphism decisive on the relevant splitting-field character line. It does not automatically give a rational isomorphism without orbit reconstruction.\par
All 364 historical K23 certificates were freshly regenerated. That closes the finite sign checks.}{Two gaps remain distinct: constructing this Chow map for the specified mechanism, and proving that the packet grammar covers all mixed-prime inputs.}

\GuideSpread{19 / The carrier begins with a six-by-six matrix}{EXACT INTEGER DATA / THE DEGREE-70 CERTIFICATE}{19_carrier70_matrix.pdf}{Rows encode the six target monomials. Columns are nonnegative exponent vectors invariant under the desired Fermat character.\par
The determinant is $-2100$. Exact arithmetic gives
\[A^{\mathsf T}a=70(3,3,1,3,2,3).\]
The congruence is what ties the monomial geometry to the desired packet.\par
The matrix was discovered through a bounded search. The recorded arithmetic certificate does not depend on the solver's floating-point tolerances.}{Do not transpose the row and column conventions. A positive matrix candidate also needs a normal, integral and well-formed geometric model. Source: [C].}

\GuideSpread{20 / Positive weights lead to a negative canonical coefficient}{EXACT INTEGER DATA / WEIGHTED DEGREE 30}{20_carrier70_weights.pdf}{The primitive weights are $W=(10,6,3,3,3,8)$. Each row monomial has weighted degree 30, because $AW=30\mathbf1$.\par
The sum of the weights is 33. The canonical coefficient is therefore $30-33=-3$.\par
The hypersurface is singular and is not quasismooth. The proof checks normality and the ambient singular-stratum condition before using adjunction.\par
It does not assume that a resolution remains Fano.}{Negative canonical intersection with a covering family of curves yields non-pseudoeffectivity on the resolution. A cited theorem then gives uniruledness.}

\GuideSpread{21 / Normality has a separate finite certificate}{EXACT COORDINATE DATA / THREE SINGULAR PLANES}{21_carrier70_singular_locus.pdf}{An entry one means that coordinate vanishes. Each row describes one affine coordinate plane in the cone's singular locus.\par
The cone has dimension five. The singular locus has dimension two, so it has codimension three.\par
The radical of the gradient ideal is
\[(y_0,y_1,y_2y_5,y_3y_4,y_3y_5).\]
Exact polynomial reductions and direct case analysis verify the locus. The hypersurface condition then supplies Cohen--Macaulayness, and Serre's criterion gives normality.}{A numerical rank or weight sum would not establish normality. The geometry has its own certificate and proof. Source: [C], Section 5.}

\GuideSpread{22 / The torus map controls the character}{EXACT INTEGER DATA / MONOMIAL COVER AT DEGREE 210}{22_carrier70_torus_matrix.pdf}{The rows of $B$ give $y_i=\prod_jx_j^{B_{ij}}$ on the coordinate torus. Negative exponents are legitimate there.\par
The identity $AB=210I$ pulls the target polynomial back to the degree-210 Fermat polynomial. The identity $B\mathbf1=7W$ makes the map projectively compatible.\par
Deck invariance, boundary vanishing and lowest-weight lifting are additional Hodge-theoretic steps.\par
The uniruled-fourfold theorem supplies cycle existence on the smooth carrier resolution.}{The result gives a special packet through an algebraic correspondence. It does not list equations for a basis of all representing surfaces.}

\GuideSpread{23 / The transfer changes the rational packet}{EXACT ORBIT DATA / DEGREE 210}{23_transfer210_orbits.pdf}{The source has 24 distinct ordered conjugates. The full-conductor target has 48. The two orbits are disjoint.\par
Every unit conjugate of both tuples has least-positive residue sum 630. Thus each is a rational middle Hodge packet.\par
Coordinate permutations are useful for multiset identities, but they do not redefine the ordered Galois orbit count.\par
The target is $b=(2,9,129,142,168,180)$.}{The positive conclusion for $b$ depends on the inherited carrier proposition and on classical algebraic Fermat operations. Source: [T].}

\GuideSpread{24 / The transfer is an exact incidence identity}{EXACT INTEGER DATA / FOURTEEN AUXILIARY TERMS}{24_transfer210_incidence.pdf}{Teal counts the multiplicities on the target side. Copper counts them on the source side. Every pair of bars agrees, including repeated residues.\par
The fourteen auxiliary generators include complementary pairs, balanced four-entry blocks and two known standard blocks. All have coefficient $+1$ or $-1$ in the signed identity.\par
The positive form has 30 entries on each side.\par
The solver found the identity. Independent integer replay checks it.}{The reported optimizer objective is not a proof that fourteen terms is globally minimal among all geometric identities.}

\GuideSpread{25 / Stabilization converts signs into geometric removal}{EXACT COUNTS / CLASSICAL CORRESPONDENCE DEPENDENCY}{25_transfer210_stabilization.pdf}{Write the identity as $b\uplus L=s\uplus R$, with $|L|=|R|=24$. Append the negated multiset $-L$.\par
The source side is algebraic by concatenation of already algebraic blocks. The target side now contains 24 complementary pairs.\par
The classical complementary-pair removal theorem returns a six-entry character.\par
The intermediate character has 54 entries and lies on a Fermat variety of dimension 52. No general Hodge theorem for that 52-fold is assumed.}{This is a geometric cancellation theorem for explicitly algebraic factors. It is not arbitrary cancellation of tensor factors or unknown motives.}

\GuideSpread{26 / Capacity measures one constructor's numerical room}{EXACT PRIMAL AND GLOBAL DUAL CERTIFICATES / SPECIFIED INPUTS}{26_capacity_values.pdf}{Capacity maximizes the total weight of invariant nonnegative exponent columns that average to $\mathbf1$. The allowed support is finite but not restricted to six independent columns.\par
The figure gives exact optima for eight specified characters. It is not a census of their degrees.\par
A value above one is a numerical opportunity. It does not certify a normal carrier. A value one excludes the strict positive-anticanonical matrix criterion.}{The degree-70 packet has capacity $11/10$. The transferred degree-210 packet has capacity one, while the transfer argument still algebraizes it.}

\GuideSpread{27 / A global obstruction can have finite premises}{EXACT INTEGER DATA PLUS A UNIFORM TAIL / THE DEGREE-210 DUAL}{27_global_dual210.pdf}{The dual vector has positive entries and total one. For invariant columns of degrees one through thirteen, exhaustive enumeration verifies $\eta\cdot v\ge1$.\par
Every column of degree at least fourteen has integer cost at least $14\cdot31=434$, exceeding the denominator 420.\par
These two statements cover every nonzero invariant column, with no search cutoff left unaddressed.\par
A separate shortest-return-path check confirms minimum cost 420.}{The finite range plus the proved tail is a global certificate. A failed bounded search without such an argument would be inconclusive.}

\GuideSpread{28 / An infinite family makes the universal matrix claim false}{SYMBOLIC MANUSCRIPT PROOF / A CONGRUENCE OBSTRUCTION}{28_fivephase_capacity_obstruction.pdf}{Assume $q>5$, $1\le b<q$ and $\gcd(b,q)=1$. Reduce the invariant-column condition modulo $q$, and cancel the unit $b$. The numerator of $y\cdot v$ is then a positive multiple of $q$.\par
Consequently $y\cdot v\ge1$ for every allowed nonzero column. Since $y\cdot\mathbf1=1$, the total normalized weight cannot exceed one.\par
The column $v=\mathbf1$ supplies the matching lower bound. Capacity is exactly one.\par
The proof holds for every stated parameter. Finite benchmarks are separate checks.}{The retracted universal claim concerns this matrix constructor. The same standard family has explicit algebraic surfaces. It gives no counterexample to Hodge.}

\GuideSpread{29 / A nonunit power blocks a direct uniruled quotient}{EXACT CHARACTER WITNESS / MODULUS 110}{29_holomorphic_power_obstruction.pdf}{The balanced character has sum $3m=330$ in every unit conjugate. Multiplying it by the nonunit eleven gives a full-support character with sum $m=110$.\par
A diagonal deck group on which the first character is trivial also kills its eleventh power. The quotient therefore carries a holomorphic top-form component.\par
Boundary control and lowest-weight lifting carry that component to any smooth projective compactification.\par
A uniruled fourfold has no holomorphic top form.}{This excludes the direct diagonal uniruled-quotient route for the witness. It leaves other carriers and non-diagonal correspondences available.}

\GuideSpread{30 / The missed family has an actual surface}{SYMBOLIC EQUATIONS / THE KNOWN FIVE-PHASE STANDARD FAMILY}{30_explicit_fivephase_surface.pdf}{Put $t_i=x_i^q$, $P=x_0x_1x_2x_3x_4$ and $z=x_5$. The three equations have degrees $q$, $2q$ and five.\par
Newton's identity reduces $p_5(t)$ to $5P^q$ on the first two equations. The relation $\kappa^q=-1/5$ then makes the Fermat polynomial vanish.\par
The proof establishes an integral complete-intersection surface of degree $10q^2$. It does not assume the surface is smooth.\par
The defining surface is Aoki's known standard construction. Its determinant and period are recomputed here, with no priority claim.}{Containment supplies an actual cycle. A nonzero primitive character projection still requires the determinant and period calculation. Source: [D], Sections 4--6.}

\GuideSpread{31 / A socle coefficient proves a nonzero projection}{EXACT EXPONENT DATA / THE DEGREE-35 EXAMPLE}{31_socle_pairing35.pdf}{At $q=7$, $b=1$, the selected determinant term has exponents $a_i-1$. The dual numerator has exponents $35-a_i-1$.\par
Each pair sums to 33, the exponent of the Fermat Jacobian socle monomial. Only this complementary determinant term survives.\par
Its coefficient is $10{,}504{,}375\kappa^6\ne0$. The complete-intersection period theorem gives the displayed nonzero period.\par
The numerical factor depends on the stated residue normalization. Nonvanishing is the invariant conclusion.}{Matching a character name is insufficient. The socle pairing supplies the actual nonzero coefficient. Source: [D], Sections 5--7.}

\GuideSpread{32 / Rationality is recovered through the whole group orbit}{EXACT DIMENSIONS / A RATIONAL ALGEBRAIC GENERATOR}{32_rational_orbit_reconstruction.pdf}{The primitive correction $\xi_\kappa=[Z_\kappa]-2q h^2$ is rational and algebraic. The diagonal group acts through algebraic automorphisms.\par
The determinant has nonzero support on the listed character lines. Multiplicity one lets their projected components generate the full rational orbit subspace.\par
At $q=7$ one ordered packet has dimension 24. The union across all first-five permutations has dimension 720.\par
These are different counts. Neither is the full Hodge dimension of the fourfold.}{The dimension $120(q-1)$ belongs to this standard-family sector with a fixed last coordinate. It is not a count of all Hodge classes.}

\GuideSpread{33 / A tangent obstruction can be local to one constructor}{EXACT LINEARIZED DATA / THE WEIL GRAPH AUDIT}{33_weil_graph_tangent_boundary.pdf}{The ambient real tangent space has dimension 18. The chosen nonsplit graph is preserved along a six-dimensional real subspace. The split control graph has dimension twelve.\par
These ranks concern those particular graphs. They do not bound deformation ranks of all possible cycles.\par
The audit also finds replacement factors along the selected mixed direction. The original factors can lose their Hodge type while the full Weil classes retain it.}{This is a negative audit of a proposed construction argument, with exact finite algebra. It proves neither nonalgebraicity nor a Hodge counterexample. Source: [W].}

\GuideSpread{34 / Results, retractions and live targets need separate rows}{STATUS MATRIX / READER EDITION 1.0}{34_status_and_targets.pdf}{The carrier and transfer arguments close specified packets at manuscript level, with fresh exact arithmetic replay. The standard surfaces recover a known algebraic family.\par
Universal positive-anticanonical matrix coverage is false, so it is retired rather than left as an open plausible lemma.\par
The K23 Chow map remains a construction target. Universal mixed-prime grammar coverage is a distinct missing implication.\par
General Hodge also requires an argument beyond these diagonal examples.}{A reproducible PASS verifies its stated finite contract. It does not turn a manuscript dependency, an external theorem or an open universal implication into a proof.}

\clearpage
{\fontsize{22}{26}\selectfont\bfseries\color{Ink}Sources and reproducibility}\par
{\sffamily\small\color{Gray}PRIMARY MATHEMATICS / PRESERVED PROJECT EVIDENCE / FRESH EXACT FIGURES}\par
\vspace{4mm}
\begin{minipage}[t]{0.48\linewidth}\small
\textbf{Primary references}\par
[1] P. Deligne, \emph{The Hodge Conjecture}, Clay Mathematics Institute problem description.\par
[2] G. da Silva Jr., \emph{Notes on the Hodge Conjecture for Fermat Varieties}, arXiv:2101.04739v2. Character decomposition and Fermat correspondences.\par
[3] D. Arapura, \emph{Varieties with very little transcendental cohomology}, arXiv:0706.2506v1, Corollary 4.2. The Conte--Murre theorem for uniruled fourfolds.\par
[4] R. Villaflor Loyola, \emph{Periods of complete intersection algebraic cycles}, arXiv:1812.03964v5, Theorems 1.1--1.2.\par
[5] N. Aoki, \emph{Some new algebraic cycles on Fermat varieties}, J. Math. Soc. Japan 39 (1987), 385--396. Classical standard cycles and operations. The inherited notes had incomplete access. The full text and cited statements are checked in this edition.\par
[6] P. Deligne, \emph{Th\'eorie de Hodge II}, Corollary 3.2.17 and Lemma 4.2.3. Lowest-weight lifting.\par
[7] S. Boucksom, J.-P. Demailly, M. P\u{a}un and T. Peternell, arXiv:math/0405285v1, Theorem 2.6. Uniruledness criterion.\par
The monograph and source audit give exact hypotheses, version checks and source boundaries.
\end{minipage}\hfill
\begin{minipage}[t]{0.48\linewidth}\small
\textbf{Project source keys}\par
[A] \emph{Projector-Resolved Hodge Calculus, Dyadic Completion}, August 2026, with source archive.\par
[C] \emph{Komorebi Hodge Carrier70}, v0.1.0.\par
[T] \emph{Komorebi Hodge Carrier Transfer}, v0.2.0.\par
[D] \emph{Komorebi Hodge Dual and Cycles}, v0.2.0.\par
[W] \emph{Hodge Weil Construction Audit}, Block 2.\par
\textbf{Figure contract}\par
All 34 figures are generated by \texttt{visual/figures.py}. The contract and data JSON identify exact inputs, checks, schematic drawings and hashes.\par
Fresh replay records are under \texttt{evidence/computation/}. Their scope includes the 24, 46 and 137 check groups, exhaustive dyadic enumeration through 128, regenerated K23 certificates and the selected Weil audit.\par
The source ZIP preserves historical failures and corrections. Figures use exact integer or rational inputs before conversion to display coordinates. No generated imagery is used.
\end{minipage}
\vfill
\colorbox{Soft}{\parbox{\dimexpr\linewidth-2\fboxsep\relax}{\small The source bundle is the reproducibility companion. The monograph supplies the full arguments. This guide supplies visual reading paths and worked certificates, with the general Hodge conjecture and the open construction obligations kept explicit.}}
\end{document}
